\chapter{How's Your Japanese Learning Going?}

% ================= PERCAKAPAN 1-5 =================
\begin{convobox}{1}
    \noindent
    \jpkana{こんにちは。}{Konnichiwa.}

    \vspace{1.5em}
    \noindent{\Large \color{articolor} \textbf{Arti:} Halo.}
\end{convobox}

\begin{convobox}{2}
    \noindent
    \jpword{にほんご}{日本語}{Nihongo} \jpkana{の}{no} 
    \jpword{べんきょう}{勉強}{benkyou}\jpkana{、}{,} 
    \jpkana{お}{o}\jpword{つか}{疲}{tsuka}\jpkana{れ}{re}\jpword{さま}{様}{sama} \jpkana{です。}{desu.}

    \vspace{1.5em}
    \noindent{\Large \color{articolor} \textbf{Arti:} Terima kasih atas kerja kerasmu belajar bahasa Jepang (Semangat belajarnya).}
\end{convobox}

\begin{convobox}{3}
    \noindent
    \jpword{すこ}{少}{suko}\jpkana{し}{shi} 
    \jpword{きゅう}{休}{kyuu}\jpkana{けいして、}{kei shite,} 
    \jpword{いっしょ}{一緒}{issho} \jpkana{に}{ni} 
    \jpkana{おしゃべりしましょう。}{oshaberi shimashou.}

    \vspace{1.5em}
    \noindent{\Large \color{articolor} \textbf{Arti:} Mari istirahat sebentar dan mengobrol bersama.}
\end{convobox}

\begin{convobox}{4}
    \noindent
    \jpword{さいきん}{最近}{Saikin}\jpkana{、}{,} 
    \jpword{にほんご}{日本語}{Nihongo} \jpkana{の}{no} 
    \jpword{べんきょう}{勉強}{benkyou} \jpkana{は}{wa} 
    \jpkana{どうですか？}{dou desu ka?}

    \vspace{1.5em}
    \noindent{\Large \color{articolor} \textbf{Arti:} Akhir-akhir ini, bagaimana dengan belajar bahasa Jepangnya?}
\end{convobox}

\begin{convobox}{5}
    \noindent
    \jpword{すこ}{少}{suko}\jpkana{し}{shi} 
    \jpword{たいへん}{大変}{taihen} \jpkana{ですが、}{desu ga,} 
    \jpword{じゅんちょう}{順調}{junchou} \jpkana{です。}{desu.}

    \vspace{1.5em}
    \noindent{\Large \color{articolor} \textbf{Arti:} Sedikit berat/susah, tapi berjalan lancar.}
\end{convobox}

% --- LESSON BUFFER 1 ---
\begin{lessonbox}{Catatan Belajar: Kata Ajaib "Otsukaresama"}
    Adik-adik, tahukah kamu kata \textbf{お疲れ様です (Otsukaresama desu)}? Ini adalah kalimat pujian super ajaib di Jepang! Artinya kurang lebih, "Terima kasih atas kerja kerasmu hari ini." \\
    
    Orang Jepang sering sekali mengucapkannya setelah selesai belajar, selesai bekerja, atau melihat temannya berusaha keras. Coba ucapkan ke teman belajarmu ya, pasti mereka akan senang!
\end{lessonbox}

% ================= PERCAKAPAN 6-12 =================
\begin{convobox}{6}
    \noindent
    \jpword{まいにち}{毎日}{Mainichi} 
    \jpkana{どれぐらいの}{dore gurai no} 
    \jpword{じかん}{時間}{jikan} 
    \jpword{べんきょう}{勉強}{benkyou} \jpkana{していますか？}{shite imasu ka?}

    \vspace{1.5em}
    \noindent{\Large \color{articolor} \textbf{Arti:} Setiap hari kamu belajar kira-kira berapa lama (waktunya)?}
\end{convobox}

\begin{convobox}{7}
    \noindent
    \jpword{いちにち}{1日}{Ichi-nichi} 
    \jpword{いちじかん}{1時間}{ichi-jikan} 
    \jpkana{ぐらい}{gurai} 
    \jpword{べんきょう}{勉強}{benkyou} \jpkana{しています。}{shite imasu.}

    \vspace{1.5em}
    \noindent{\Large \color{articolor} \textbf{Arti:} Aku belajar kira-kira 1 jam sehari.}
\end{convobox}

\begin{convobox}{8}
    \noindent
    \jpkana{いつ}{Itsu} 
    \jpword{べんきょう}{勉強}{benkyou} \jpkana{していますか？}{shite imasu ka?}

    \vspace{1.5em}
    \noindent{\Large \color{articolor} \textbf{Arti:} Kapan kamu belajar?}
\end{convobox}

\begin{convobox}{9}
    \noindent
    \jpword{あさ}{朝}{Asa} \jpkana{ですか？}{desu ka?} 
    \jpword{よる}{夜}{Yoru} \jpkana{ですか？}{desu ka?}

    \vspace{1.5em}
    \noindent{\Large \color{articolor} \textbf{Arti:} Apakah pagi hari? Apakah malam hari?}
\end{convobox}

\begin{convobox}{10}
    \noindent
    \jpword{あさ}{朝}{Asa}\jpkana{ごはんを}{gohan o} 
    \jpword{た}{食}{ta}\jpkana{べた}{beta} 
    \jpword{あと}{後}{ato} \jpkana{と、}{to,}

    \vspace{1.5em}
    \noindent{\Large \color{articolor} \textbf{Arti:} Setelah makan sarapan (pagi), dan}
\end{convobox}

\begin{convobox}{11}
    \noindent
    \jpword{よる}{夜}{Yoru} 
    \jpword{ね}{寝}{ne}\jpkana{る}{ru} 
    \jpword{まえ}{前}{mae} \jpkana{に}{ni} 
    \jpword{べんきょう}{勉強}{benkyou} \jpkana{しています。}{shite imasu.}

    \vspace{1.5em}
    \noindent{\Large \color{articolor} \textbf{Arti:} sebelum tidur di malam hari, aku belajar.}
\end{convobox}

\begin{convobox}{12}
    \noindent
    \jpkana{どうやって}{Dou yatte} 
    \jpword{にほんご}{日本語}{Nihongo} \jpkana{を}{o} 
    \jpword{べんきょう}{勉強}{benkyou} \jpkana{を}{o} 
    \jpkana{していますか？}{shite imasu ka?}

    \vspace{1.5em}
    \noindent{\Large \color{articolor} \textbf{Arti:} Bagaimana caramu belajar bahasa Jepang?}
\end{convobox}

% --- LESSON BUFFER 2 ---
\begin{lessonbox}{Catatan Belajar: Sebelum dan Sesudah!}
    Waktu belajar setiap orang beda-beda ya! Ada yang pagi, ada yang malam. Nah, kalau kita mau bilang "Sesudah" dan "Sebelum", kita bisa pakai rumus ini:
    \begin{itemize}
        \item \textbf{〜た後 (〜ta ato)} = Setelah (melakukan sesuatu). Contoh: \textit{Tabeta ato} (Setelah makan).
        \item \textbf{〜る前 (〜ru mae)} = Sebelum (melakukan sesuatu). Contoh: \textit{Neru mae} (Sebelum tidur).
    \end{itemize}
    Coba deh ingat-ingat, kegiatan apa yang biasa kamu lakukan "Neru mae" (sebelum tidur)? Sikat gigi? Baca buku? Atau belajar bahasa Jepang?
\end{lessonbox}

% ================= PERCAKAPAN 13-18 =================
\begin{convobox}{13}
    \noindent
    \jpkana{お}{O}\jpword{き}{気}{ki} \jpkana{に}{ni} \jpword{い}{入}{i}\jpkana{りの}{ri no} 
    \jpword{きょうかしょ}{教科書}{kyoukasho} \jpkana{や、}{ya,} 
    \jpkana{Webサイト}{Web saito} \jpkana{は}{wa} 
    \jpkana{ありますか？}{arimasu ka?}

    \vspace{1.5em}
    \noindent{\Large \color{articolor} \textbf{Arti:} Apakah ada buku pelajaran atau website favorit?}
\end{convobox}

\begin{convobox}{14}
    \noindent
    \jpword{きょうかしょ}{教科書}{Kyoukasho} \jpkana{と}{to} 
    \jpkana{YouTube}{YouTube} \jpkana{を}{o} 
    \jpword{つか}{使}{tsuka}\jpkana{って}{tte} 
    \jpword{べんきょう}{勉強}{benkyou} \jpkana{しています。}{shite imasu.}

    \vspace{1.5em}
    \noindent{\Large \color{articolor} \textbf{Arti:} Aku belajar menggunakan buku pelajaran dan YouTube.}
\end{convobox}

\begin{convobox}{15}
    \noindent
    \jpkana{YouTube}{YouTube} \jpkana{は、}{wa,} 
    \jpkana{たなかさんの}{Tanaka-san no} 
    \jpword{どうが}{動画}{douga} \jpkana{が}{ga} 
    \jpkana{お}{o}\jpword{き}{気}{ki} \jpkana{に}{ni} \jpword{い}{入}{i}\jpkana{りです。}{ri desu.}

    \vspace{1.5em}
    \noindent{\Large \color{articolor} \textbf{Arti:} Di YouTube, video dari Tanaka-san adalah favoritku.}
\end{convobox}

\begin{convobox}{16}
    \noindent
    \jpword{さいきん}{最近}{Saikin} 
    \jpword{おぼ}{覚}{obo}\jpkana{えた}{eta} 
    \jpword{ことば}{言葉}{kotoba} \jpkana{は}{wa} 
    \jpword{なに}{何}{nani}\jpkana{か}{ka} 
    \jpkana{ありますか？}{arimasu ka?}

    \vspace{1.5em}
    \noindent{\Large \color{articolor} \textbf{Arti:} Apakah ada kosakata yang baru-baru ini kamu hafal/ingat?}
\end{convobox}

\begin{convobox}{17}
    \noindent
    \jpkana{「チョキチョキ」という}{"Choki-choki" to iu} 
    \jpword{ことば}{言葉}{kotoba} \jpkana{を}{o} 
    \jpword{おぼ}{覚}{obo}\jpkana{えました。}{emashita.}

    \vspace{1.5em}
    \noindent{\Large \color{articolor} \textbf{Arti:} Aku menghafal kata "Choki-choki".}
\end{convobox}

\begin{convobox}{18}
    \noindent
    \jpkana{はさみで}{Hasami de} 
    \jpword{なに}{何}{nani}\jpkana{か}{ka} \jpkana{を}{o} 
    \jpword{き}{切}{ki}\jpkana{る}{ru} 
    \jpword{おと}{音}{oto} \jpkana{を}{o} 
    \jpword{あらわ}{表}{arawa}\jpkana{す}{su} 
    \jpkana{オノマトペです。}{onomatope desu.}

    \vspace{1.5em}
    \noindent{\Large \color{articolor} \textbf{Arti:} Itu adalah Onomatope yang menggambarkan suara memotong sesuatu dengan gunting.}
\end{convobox}

% --- LESSON BUFFER 3 ---
\begin{lessonbox}{Catatan Belajar: Kata Benda Kesukaan \& Onomatope!}
    Kata \textbf{お気に入り (O-ki ni iri)} artinya "Favorit" atau kesukaan. Jadi kalau kamu mau bilang "Buku favorit", tinggal tambahkan \textbf{no} jadi \textit{O-ki ni iri no hon}! \\
    
    Nah, terus apa itu \textbf{オノマトペ (Onomatope)}? Itu adalah kata tiruan bunyi! Di Indonesia kita punya kata "Dor!" untuk suara tembakan. Di Jepang, suara gunting memotong kertas itu bunyinya adalah \textbf{チョキチョキ (Choki-choki)}! Lucu banget, ya?
\end{lessonbox}

% ================= PERCAKAPAN 19-24 =================
\begin{convobox}{19}
    \noindent
    \jpkana{ふだん}{Fudan} 
    \jpword{ひとり}{一人}{hitori} \jpkana{で}{de} 
    \jpword{べんきょう}{勉強}{benkyou} \jpkana{していますか？}{shite imasu ka?}

    \vspace{1.5em}
    \noindent{\Large \color{articolor} \textbf{Arti:} Biasanya kamu belajar sendirian?}
\end{convobox}

\begin{convobox}{20}
    \noindent
    \jpkana{それとも、}{Sore tomo,} 
    \jpword{せんせい}{先生}{sensei} \jpkana{や}{ya} 
    \jpword{ともだち}{友達}{tomodachi} \jpkana{と}{to} 
    \jpword{いっしょ}{一緒}{issho} \jpkana{に}{ni} 
    \jpword{べんきょう}{勉強}{benkyou} \jpkana{していますか？}{shite imasu ka?}

    \vspace{1.5em}
    \noindent{\Large \color{articolor} \textbf{Arti:} Atau, kamu belajar bersama guru dan teman-teman?}
\end{convobox}

\begin{convobox}{21}
    \noindent
    \jpword{ひとり}{一人}{Hitori} \jpkana{で}{de} 
    \jpword{べんきょう}{勉強}{benkyou} \jpkana{しています。}{shite imasu.}

    \vspace{1.5em}
    \noindent{\Large \color{articolor} \textbf{Arti:} Aku belajar sendirian.}
\end{convobox}

\begin{convobox}{22}
    \noindent
    \jpword{べんきょう}{勉強}{Benkyou} \jpkana{の}{no} 
    \jpkana{やる}{yaru}\jpword{き}{気}{ki} \jpkana{が}{ga} 
    \jpword{で}{出}{de}\jpkana{ない}{nai} 
    \jpword{とき}{時}{toki} \jpkana{は}{wa} 
    \jpkana{どうしていますか？}{dou shite imasu ka?}

    \vspace{1.5em}
    \noindent{\Large \color{articolor} \textbf{Arti:} Apa yang kamu lakukan saat sedang tidak ada motivasi belajar?}
\end{convobox}

\begin{convobox}{23}
    \noindent
    \jpword{べんきょう}{勉強}{Benkyou} \jpkana{を}{o} 
    \jpword{すこ}{少}{suko}\jpkana{し}{shi} 
    \jpword{やす}{休}{yasu}\jpkana{んで、}{nde,} 
    \jpword{さんぽ}{散歩}{sanpo} \jpkana{に}{ni} 
    \jpword{い}{行}{i}\jpkana{ったり、}{ttari,}

    \vspace{1.5em}
    \noindent{\Large \color{articolor} \textbf{Arti:} Aku istirahat belajar sebentar, lalu pergi jalan-jalan (santai), atau}
\end{convobox}

\begin{convobox}{24}
    \noindent
    \jpkana{ジムに}{Jimu ni} 
    \jpword{い}{行}{i}\jpkana{ったりしています。}{ttari shite imasu.}

    \vspace{1.5em}
    \noindent{\Large \color{articolor} \textbf{Arti:} pergi ke tempat gym.}
\end{convobox}

% --- LESSON BUFFER 4 ---
\begin{lessonbox}{Catatan Belajar: Kehilangan Semangat?}
    Ada kalanya kita merasa malas atau tidak semangat belajar. Dalam bahasa Jepang, motivasi atau kemauan disebut \textbf{やる気 (Yaruki)}. Kalau motivasinya tidak muncul, kita bilang \textit{Yaruki ga denai}. \\
    
    Kalau sedang begitu, tidak apa-apa beristirahat sebentar! Di percakapan ini, kita pakai tata bahasa \textbf{〜たり、〜たり (〜tari, 〜tari)} lagi lho, seperti yang kita pelajari di Bab 3. Gunanya untuk menyebutkan kegiatan-kegiatan acak yang kita lakukan saat sedang istirahat.
\end{lessonbox}

% ================= PERCAKAPAN 25-37 =================
\begin{convobox}{25}
    \noindent
    \jpword{にほんご}{日本語}{Nihongo} \jpkana{の}{no} 
    \jpword{うた}{歌}{uta} \jpkana{を}{o} 
    \jpkana{よく}{yoku} 
    \jpword{き}{聴}{ki}\jpkana{きますか？}{kimasu ka?}

    \vspace{1.5em}
    \noindent{\Large \color{articolor} \textbf{Arti:} Apakah kamu sering mendengarkan lagu bahasa Jepang?}
\end{convobox}

\begin{convobox}{26}
    \noindent
    \jpkana{おすすめの}{Osusume no} 
    \jpword{うた}{歌}{uta} \jpkana{が}{ga} 
    \jpkana{あったら、}{attara,} 
    \jpword{おし}{教}{oshi}\jpkana{えてください。}{ete kudasai.}

    \vspace{1.5em}
    \noindent{\Large \color{articolor} \textbf{Arti:} Kalau ada lagu rekomendasi, tolong beritahu ya.}
\end{convobox}

\begin{convobox}{27}
    \noindent
    \jpword{にほんご}{日本語}{Nihongo} \jpkana{の}{no} 
    \jpword{うた}{歌}{uta} \jpkana{は}{wa} 
    \jpkana{あまり}{amari} 
    \jpword{き}{聴}{ki}\jpkana{きません。}{kimasen.}

    \vspace{1.5em}
    \noindent{\Large \color{articolor} \textbf{Arti:} Aku tidak terlalu sering mendengarkan lagu bahasa Jepang.}
\end{convobox}

\begin{convobox}{28}
    \noindent
    \jpkana{おすすめを}{Osusume o} 
    \jpword{し}{知}{shi}\jpkana{りたいです。}{ritai desu.}

    \vspace{1.5em}
    \noindent{\Large \color{articolor} \textbf{Arti:} Aku yang ingin tahu rekomendasinya.}
\end{convobox}

\begin{convobox}{29}
    \noindent
    \jpword{にほんご}{日本語}{Nihongo} \jpkana{が}{ga} 
    \jpword{いま}{今}{ima} \jpkana{よりもっと}{yori motto} 
    \jpword{じょうず}{上手}{jouzu} \jpkana{になったら、}{ni nattara,}

    \vspace{1.5em}
    \noindent{\Large \color{articolor} \textbf{Arti:} Kalau bahasa Jepangmu nanti sudah menjadi lebih jago dari sekarang,}
\end{convobox}

\begin{convobox}{30}
    \noindent
    \jpword{なに}{何}{Nani} \jpkana{を}{o} 
    \jpkana{したいですか？}{shitai desu ka?}

    \vspace{1.5em}
    \noindent{\Large \color{articolor} \textbf{Arti:} apa yang ingin kamu lakukan?}
\end{convobox}

\begin{convobox}{31}
    \noindent
    \jpword{にほん}{日本}{Nihon}\jpword{りょこう}{旅行}{ryokou} \jpkana{をして、}{o shite,} 
    \jpword{げんち}{現地}{genchi} \jpkana{の}{no} 
    \jpword{ひと}{人}{hito} \jpkana{と}{to} 
    \jpkana{たくさん}{takusan} 
    \jpword{はな}{話}{hana}\jpkana{したいです。}{shitai desu.}

    \vspace{1.5em}
    \noindent{\Large \color{articolor} \textbf{Arti:} Aku ingin jalan-jalan (wisata) di Jepang, dan banyak berbicara dengan penduduk lokal.}
\end{convobox}

\begin{convobox}{32}
    \noindent
    \jpkana{とてもすてきな}{Totemo suteki na} 
    \jpword{もくひょう}{目標}{mokuhyou} \jpkana{ですね。}{desu ne.}

    \vspace{1.5em}
    \noindent{\Large \color{articolor} \textbf{Arti:} Itu adalah tujuan/target yang sangat indah ya.}
\end{convobox}

\begin{convobox}{33}
    \noindent
    \jpword{きょう}{今日}{Kyou} \jpkana{は}{wa} 
    \jpkana{お}{o}\jpword{はなし}{話}{hanashi} \jpkana{できて}{dekite} 
    \jpword{たの}{楽}{tano}\jpkana{しかったです。}{shikatta desu.}

    \vspace{1.5em}
    \noindent{\Large \color{articolor} \textbf{Arti:} Hari ini aku senang sekali bisa mengobrol.}
\end{convobox}

\begin{convobox}{34}
    \noindent
    \jpword{べんきょう}{勉強}{Benkyou} \jpkana{が}{ga} 
    \jpkana{うまくいかない}{umaku ikanai} 
    \jpword{とき}{時}{toki} \jpkana{や、}{ya,}

    \vspace{1.5em}
    \noindent{\Large \color{articolor} \textbf{Arti:} Saat belajarnya sedang tidak berjalan lancar, atau}
\end{convobox}

\begin{convobox}{35}
    \noindent
    \jpkana{ちょっと}{Chotto} 
    \jpword{きゅう}{休}{kyuu}\jpkana{けいしたい}{kei shitai} 
    \jpword{とき}{時}{toki} \jpkana{は、}{wa,}

    \vspace{1.5em}
    \noindent{\Large \color{articolor} \textbf{Arti:} saat ingin sedikit beristirahat,}
\end{convobox}

\begin{convobox}{36}
    \noindent
    \jpkana{また}{Mata} \jpkana{いつでも}{itsu demo} 
    \jpkana{ここに}{koko ni} 
    \jpword{き}{来}{ki}\jpkana{てくださいね。}{te kudasai ne.}

    \vspace{1.5em}
    \noindent{\Large \color{articolor} \textbf{Arti:} silakan datang ke sini kapan saja ya.}
\end{convobox}

\begin{convobox}{37}
    \noindent
    \jpkana{それでは。}{Sore dewa.}

    \vspace{1.5em}
    \noindent{\Large \color{articolor} \textbf{Arti:} Kalau begitu, (sampai jumpa).}
\end{convobox}

% --- SUMMARY BOX ---
\vspace{2em}
\begin{tcolorbox}[
    colback=blue!5!white,
    colframe=blue!80!black,
    boxrule=3pt,
    arc=5mm,
    title={\huge\bfseries \sffamily 🌟 Rangkuman Bab 6 🌟},
    fonttitle=\color{white},
    attach boxed title to top center={yshift=-4mm},
    boxed title style={colback=blue!80!black, rounded corners},
    left=5mm, right=5mm, top=7mm, bottom=5mm
]
\Large \textbf{Otsukaresama desu! Kamu telah menyelesaikan Bab 6!} Yuk kita cek apa saja yang sudah kita pelajari:
\begin{itemize}
    \item \textbf{Kalimat Pujian:} \textit{Otsukaresama desu} (Terima kasih atas kerja kerasmu).
    \item \textbf{Kata Keterangan Waktu:} \textit{\dots ta ato} (Sesudah) dan \textit{\dots ru mae ni} (Sebelum).
    \item \textbf{Kesukaan / Favorit:} \textit{O-ki ni iri}.
    \item \textbf{Tiruan Bunyi:} \textit{Onomatope}. Contoh: \textit{Choki-choki} (Suara gunting).
    \item \textbf{Motivasi Belajar:} \textit{Yaruki}. Kalau lagi malas, bilang saja \textit{Yaruki ga denai}.
\end{itemize}
Ingat pesan di akhir percakapan ya! Kalau kamu capek atau belajarnya terasa sulit, tidak apa-apa untuk istirahat sejenak. Kamu punya tujuan yang hebat, jadi pelan-pelan saja yang penting pasti!
\end{tcolorbox}

\newpage